% Options for packages loaded elsewhere
\PassOptionsToPackage{unicode}{hyperref}
\PassOptionsToPackage{hyphens}{url}
\documentclass[
  UTF8,
  a4paper,
  12pt]{ctexart}
\usepackage{xcolor}
\usepackage[margin=2.5cm]{geometry}
\usepackage{amsmath,amssymb}
\setcounter{secnumdepth}{-\maxdimen} % remove section numbering
\usepackage{iftex}
\ifPDFTeX
  \usepackage[T1]{fontenc}
  \usepackage[utf8]{inputenc}
  \usepackage{textcomp} % provide euro and other symbols
\else % if luatex or xetex
  \usepackage{unicode-math} % this also loads fontspec
  \defaultfontfeatures{Scale=MatchLowercase}
  \defaultfontfeatures[\rmfamily]{Ligatures=TeX,Scale=1}
\fi
\usepackage{lmodern}
\ifPDFTeX\else
  % xetex/luatex font selection
\fi
% Use upquote if available, for straight quotes in verbatim environments
\IfFileExists{upquote.sty}{\usepackage{upquote}}{}
\IfFileExists{microtype.sty}{% use microtype if available
  \usepackage[]{microtype}
  \UseMicrotypeSet[protrusion]{basicmath} % disable protrusion for tt fonts
}{}
\makeatletter
\@ifundefined{KOMAClassName}{% if non-KOMA class
  \IfFileExists{parskip.sty}{%
    \usepackage{parskip}
  }{% else
    \setlength{\parindent}{0pt}
    \setlength{\parskip}{6pt plus 2pt minus 1pt}}
}{% if KOMA class
  \KOMAoptions{parskip=half}}
\makeatother
\usepackage{color}
\usepackage{fancyvrb}
\newcommand{\VerbBar}{|}
\newcommand{\VERB}{\Verb[commandchars=\\\{\}]}
\DefineVerbatimEnvironment{Highlighting}{Verbatim}{commandchars=\\\{\}}
% Add ',fontsize=\small' for more characters per line
\newenvironment{Shaded}{}{}
\newcommand{\AlertTok}[1]{\textcolor[rgb]{1.00,0.00,0.00}{\textbf{#1}}}
\newcommand{\AnnotationTok}[1]{\textcolor[rgb]{0.38,0.63,0.69}{\textbf{\textit{#1}}}}
\newcommand{\AttributeTok}[1]{\textcolor[rgb]{0.49,0.56,0.16}{#1}}
\newcommand{\BaseNTok}[1]{\textcolor[rgb]{0.25,0.63,0.44}{#1}}
\newcommand{\BuiltInTok}[1]{\textcolor[rgb]{0.00,0.50,0.00}{#1}}
\newcommand{\CharTok}[1]{\textcolor[rgb]{0.25,0.44,0.63}{#1}}
\newcommand{\CommentTok}[1]{\textcolor[rgb]{0.38,0.63,0.69}{\textit{#1}}}
\newcommand{\CommentVarTok}[1]{\textcolor[rgb]{0.38,0.63,0.69}{\textbf{\textit{#1}}}}
\newcommand{\ConstantTok}[1]{\textcolor[rgb]{0.53,0.00,0.00}{#1}}
\newcommand{\ControlFlowTok}[1]{\textcolor[rgb]{0.00,0.44,0.13}{\textbf{#1}}}
\newcommand{\DataTypeTok}[1]{\textcolor[rgb]{0.56,0.13,0.00}{#1}}
\newcommand{\DecValTok}[1]{\textcolor[rgb]{0.25,0.63,0.44}{#1}}
\newcommand{\DocumentationTok}[1]{\textcolor[rgb]{0.73,0.13,0.13}{\textit{#1}}}
\newcommand{\ErrorTok}[1]{\textcolor[rgb]{1.00,0.00,0.00}{\textbf{#1}}}
\newcommand{\ExtensionTok}[1]{#1}
\newcommand{\FloatTok}[1]{\textcolor[rgb]{0.25,0.63,0.44}{#1}}
\newcommand{\FunctionTok}[1]{\textcolor[rgb]{0.02,0.16,0.49}{#1}}
\newcommand{\ImportTok}[1]{\textcolor[rgb]{0.00,0.50,0.00}{\textbf{#1}}}
\newcommand{\InformationTok}[1]{\textcolor[rgb]{0.38,0.63,0.69}{\textbf{\textit{#1}}}}
\newcommand{\KeywordTok}[1]{\textcolor[rgb]{0.00,0.44,0.13}{\textbf{#1}}}
\newcommand{\NormalTok}[1]{#1}
\newcommand{\OperatorTok}[1]{\textcolor[rgb]{0.40,0.40,0.40}{#1}}
\newcommand{\OtherTok}[1]{\textcolor[rgb]{0.00,0.44,0.13}{#1}}
\newcommand{\PreprocessorTok}[1]{\textcolor[rgb]{0.74,0.48,0.00}{#1}}
\newcommand{\RegionMarkerTok}[1]{#1}
\newcommand{\SpecialCharTok}[1]{\textcolor[rgb]{0.25,0.44,0.63}{#1}}
\newcommand{\SpecialStringTok}[1]{\textcolor[rgb]{0.73,0.40,0.53}{#1}}
\newcommand{\StringTok}[1]{\textcolor[rgb]{0.25,0.44,0.63}{#1}}
\newcommand{\VariableTok}[1]{\textcolor[rgb]{0.10,0.09,0.49}{#1}}
\newcommand{\VerbatimStringTok}[1]{\textcolor[rgb]{0.25,0.44,0.63}{#1}}
\newcommand{\WarningTok}[1]{\textcolor[rgb]{0.38,0.63,0.69}{\textbf{\textit{#1}}}}
\usepackage{longtable,booktabs,array}
\newcounter{none} % for unnumbered tables
\usepackage{calc} % for calculating minipage widths
% Correct order of tables after \paragraph or \subparagraph
\usepackage{etoolbox}
\makeatletter
\patchcmd\longtable{\par}{\if@noskipsec\mbox{}\fi\par}{}{}
\makeatother
% Allow footnotes in longtable head/foot
\IfFileExists{footnotehyper.sty}{\usepackage{footnotehyper}}{\usepackage{footnote}}
\makesavenoteenv{longtable}
\setlength{\emergencystretch}{3em} % prevent overfull lines
\providecommand{\tightlist}{%
  \setlength{\itemsep}{0pt}\setlength{\parskip}{0pt}}
\usepackage{bookmark}
\IfFileExists{xurl.sty}{\usepackage{xurl}}{} % add URL line breaks if available
\urlstyle{same}
\hypersetup{
  pdftitle={Optimal Resizable Arrays --- 统一 Review},
  hidelinks,
  pdfcreator={LaTeX via pandoc}}

\title{Optimal Resizable Arrays --- 统一 Review}
\author{}
\date{}

\begin{document}
\maketitle

\begin{quote}
论文：Optimal Resizable Arrays (Tarjan \& Zwick, 2024)
仓库：\href{https://github.com/your-org/algorithm-pa6-optimal-resizable-arrays}{algorithm-pa6-optimal-resizable-arrays}
\end{quote}

\begin{center}\rule{0.5\linewidth}{0.5pt}\end{center}

\section{第1章 摘要}\label{ux7b2c1ux7ae0-ux6458ux8981}

动态数组（Resizable Array）是一种支持自动调整大小的数据结构。

它具有以下特点：

\begin{itemize}
\tightlist
\item
  数组长度可以随着元素数量变化动态增加或减少。
\item
  支持通过索引访问元素。
\item
  随机访问时间保持为 \(O(1)\)。
\end{itemize}

Tarjan 和 Zwick 在论文《Optimal Resizable
Arrays》中研究了动态数组中的空间与时间权衡问题。

论文的核心思想：

\begin{itemize}
\tightlist
\item
  将空间划分为存储空间（storage space）和临时空间（temporary space）。
\item
  利用调整过程中的额外临时空间，降低长期存储空间需求。
\end{itemize}

论文证明，对于任意整数 \(r \geq 2\)，可以实现：

{\def\LTcaptype{none} % do not increment counter
\begin{longtable}[]{@{}ll@{}}
\toprule\noalign{}
项目 & 复杂度 \\
\midrule\noalign{}
\endhead
\bottomrule\noalign{}
\endlastfoot
存储空间 & \(N + O(N^{1/r})\) \\
临时空间 & \(N + O(N^{1-1/r})\) \\
随机访问 & \(O(1)\) \\
扩容/缩容 & \(O(r)\) 均摊时间 \\
\end{longtable}
}

该论文的主要贡献包括：

\begin{enumerate}
\def\labelenumi{\arabic{enumi}.}
\tightlist
\item
  重新定义动态数组优化问题的研究方向。
\item
  在保持高效操作的同时进一步降低存储空间。
\item
  给出了动态数组空间和时间复杂度之间更加精细的理论权衡。
\end{enumerate}

从理论角度看，该工作证明了动态数组仍然存在进一步优化空间。

从工程角度看，该论文展示了一种典型设计思想：

通过增加调整过程中的复杂度，换取长期运行时更高的空间利用率。

因此，本文认为《Optimal Resizable Arrays》的价值主要体现在两个方面：

\begin{itemize}
\tightlist
\item
  推进了动态数组空间优化的理论边界。
\item
  为其他动态数据结构的空间设计提供了新的优化思路。
\end{itemize}

\begin{center}\rule{0.5\linewidth}{0.5pt}\end{center}

\section{第2章 引言}\label{ux7b2c2ux7ae0-ux5f15ux8a00}

动态数组是现代软件系统中的基础数据结构。

例如：

\begin{itemize}
\tightlist
\item
  C++ \texttt{vector}。
\item
  Java \texttt{ArrayList}。
\item
  Python \texttt{list}。
\end{itemize}

这些实现通常采用类似倍增策略（doubling technique）。

当数组空间不足时：

\begin{enumerate}
\def\labelenumi{\arabic{enumi}.}
\tightlist
\item
  创建一个更大的连续存储空间。
\item
  将已有元素复制到新空间。
\item
  释放旧空间。
\end{enumerate}

这种方法具有较好的时间效率：

\begin{itemize}
\tightlist
\item
  单次访问时间为 \(O(1)\)。
\item
  扩容操作均摊时间为 \(O(1)\)。
\end{itemize}

但其空间利用率有限：

\begin{itemize}
\tightlist
\item
  扩容后会保留大量未使用空间。
\item
  最坏情况下接近一半容量处于空闲状态。
\end{itemize}

假设当前数组包含 \(N\) 个元素，容量可能达到 \(2N\)：

\begin{Shaded}
\begin{Highlighting}[]
\NormalTok{|{-}{-}{-}{-}{-}{-}{-}{-}{-}{-}{-} 已使用 {-}{-}{-}{-}{-}{-}{-}{-}{-}{-}{-}|{-}{-}{-}{-}{-}{-}{-}{-}{-}{-}{-} 空闲 {-}{-}{-}{-}{-}{-}{-}{-}{-}{-}{-}|}

\NormalTok{             N                     N}
\end{Highlighting}
\end{Shaded}

其中一部分空间长期处于未使用状态。

因此，传统动态数组存在一个核心问题：

如何在保持高效访问和调整大小操作的同时，进一步降低存储空间。

Tarjan 和 Zwick 在论文《Optimal Resizable Arrays》中研究了这一问题。

论文的核心思想：

\begin{itemize}
\tightlist
\item
  将空间划分为存储空间（storage space）和临时空间（temporary space）。
\item
  利用调整过程中的额外临时空间，降低长期存储空间需求。
\end{itemize}

论文关注动态数组中的空间和时间权衡。

主要目标：

\begin{enumerate}
\def\labelenumi{\arabic{enumi}.}
\tightlist
\item
  保持随机访问时间为 \(O(1)\)。
\item
  保持 Grow 和 Shrink 操作具有较低均摊时间。
\item
  将长期存储空间降低到接近实际元素数量 \(N\)。
\end{enumerate}

Tarjan 和 Zwick 给出的结果如下：

{\def\LTcaptype{none} % do not increment counter
\begin{longtable}[]{@{}ll@{}}
\toprule\noalign{}
项目 & 复杂度 \\
\midrule\noalign{}
\endhead
\bottomrule\noalign{}
\endlastfoot
Storage Space & \(N+O(N^{1/r})\) \\
Temporary Space & \(N+O(N^{1-1/r})\) \\
Access & \(O(1)\) \\
Grow/Shrink & \(O(r)\) 均摊时间 \\
\end{longtable}
}

其中参数 \(r\) 用于控制空间和时间之间的权衡：

\begin{itemize}
\tightlist
\item
  增大 \(r\) 可以进一步减少存储空间。
\item
  同时会增加调整过程的时间成本。
\end{itemize}

该论文的重要性包括：

\begin{enumerate}
\def\labelenumi{\arabic{enumi}.}
\tightlist
\item
  降低动态数组长期存储空间。
\item
  给出空间和时间复杂度之间新的理论界限。
\item
  展示通过增加调整成本换取空间优化的方法。
\end{enumerate}

该研究说明，即使是经典数据结构，也存在进一步优化空间。

本文以下章节组织如下：第3章给出预备知识，包括动态数组的形式化定义和空间模型；第4章综述相关工作和研究脉络；第5章分析下界理论；第6章展开上界构造；第7章通过
Growth Game
证明紧致性；第8章进行技术讨论与评价；第9章讨论局限性与开放问题；第10章总结全文。

\begin{center}\rule{0.5\linewidth}{0.5pt}\end{center}

\section{第3章 预备知识}\label{ux7b2c3ux7ae0-ux9884ux5907ux77e5ux8bc6}

\subsection{3.1 Resizable Array
定义}\label{resizable-array-ux5b9aux4e49}

Resizable Array 是一种支持动态调整大小的数组结构。

它需要同时满足：

\begin{itemize}
\tightlist
\item
  数组长度可以动态增加或减少。
\item
  元素可以通过索引直接访问。
\item
  访问操作保持较低时间复杂度。
\end{itemize}

论文研究的动态数组支持以下基本操作：

{\def\LTcaptype{none} % do not increment counter
\begin{longtable}[]{@{}ll@{}}
\toprule\noalign{}
操作 & 功能 \\
\midrule\noalign{}
\endhead
\bottomrule\noalign{}
\endlastfoot
\texttt{Array()} & 创建空数组 \\
\texttt{Length()} & 返回当前数组长度 \\
\texttt{Get(i)} & 返回索引 \(i\) 位置元素 \\
\texttt{Set(i,x)} & 修改索引 \(i\) 位置元素 \\
\texttt{Grow(x)} & 在数组末尾增加元素 \\
\texttt{Shrink()} & 删除数组末尾元素 \\
\end{longtable}
}

动态数组的核心目标是在保持访问效率的同时，降低扩容和缩容产生的额外空间开销。

\subsection{3.2 空间模型}\label{ux7a7aux95f4ux6a21ux578b}

论文将动态数组占用的空间分为两类：

\subsubsection{Storage Space}\label{storage-space}

Storage Space 指数组稳定运行状态下保存元素所需要的空间。

对于大小为 \(N\) 的数组：

\[
N+s(N)
\]

其中：

\begin{itemize}
\tightlist
\item
  \(N\) 表示实际存储元素数量。
\item
  \(s(N)\) 表示额外存储空间。
\end{itemize}

\subsubsection{Temporary Space}\label{temporary-space}

Temporary Space 指扩容或缩容过程中短暂使用的额外空间。

对于大小为 \(N\) 的数组：

\[
N+t(N)
\]

该模型允许调整过程中使用更多临时空间，从而降低长期存储空间需求。

\begin{center}\rule{0.5\linewidth}{0.5pt}\end{center}

\section{第4章 相关工作}\label{ux7b2c4ux7ae0-ux76f8ux5173ux5de5ux4f5c}

动态数组的发展经历了从简单倍增方案到空间优化结构的发展过程。

整体发展路线：

\begin{Shaded}
\begin{Highlighting}[]
\NormalTok{Doubling}
\NormalTok{    ↓}
\NormalTok{Sitarski HAT}
\NormalTok{    ↓}
\NormalTok{Brodnik 等人的空间下界研究}
\NormalTok{    ↓}
\NormalTok{Tarjan{-}Zwick Optimal Resizable Arrays}
\end{Highlighting}
\end{Shaded}

\subsection{4.1
传统倍增方法（Doubling）}\label{ux4f20ux7edfux500dux589eux65b9ux6cd5doubling}

传统动态数组通常维护一个连续存储空间。

当容量不足时：

\begin{itemize}
\tightlist
\item
  创建更大的数组。
\item
  复制已有元素。
\item
  释放旧数组。
\end{itemize}

特点：

{\def\LTcaptype{none} % do not increment counter
\begin{longtable}[]{@{}ll@{}}
\toprule\noalign{}
项目 & 表现 \\
\midrule\noalign{}
\endhead
\bottomrule\noalign{}
\endlastfoot
访问时间 & \(O(1)\) \\
扩容时间 & \(O(1)\) 均摊 \\
实现复杂度 & 低 \\
\end{longtable}
}

优势：

\begin{itemize}
\tightlist
\item
  实现简单。
\item
  工程应用广泛。
\end{itemize}

限制：

\begin{itemize}
\tightlist
\item
  存储空间可能达到 \(2N\)。
\item
  存在较大的空闲空间。
\end{itemize}

\subsection{4.2 Sitarski 的 HAT Data
Structure}\label{sitarski-ux7684-hat-data-structure}

Sitarski 提出了 HAT（Hierarchical Array Tree）结构。

核心思想：

\begin{itemize}
\tightlist
\item
  将数组拆分成多个存储块。
\item
  使用额外索引结构管理这些块。
\end{itemize}

特点：

{\def\LTcaptype{none} % do not increment counter
\begin{longtable}[]{@{}ll@{}}
\toprule\noalign{}
项目 & 表现 \\
\midrule\noalign{}
\endhead
\bottomrule\noalign{}
\endlastfoot
存储空间 & \(N+O(\sqrt{N})\) \\
操作时间 & \(O(1)\) \\
\end{longtable}
}

相比传统倍增方法，HAT 降低了空间浪费。

\subsection{4.3 Brodnik
等人的研究}\label{brodnik-ux7b49ux4ebaux7684ux7814ux7a76}

Brodnik 等人进一步研究动态数组的空间下界问题。

他们证明：

在限制额外临时空间的情况下，支持动态调整大小和随机访问的数据结构需要至少：

\[
N+\Omega(\sqrt{N})
\]

的额外空间。

这一结果说明：

\begin{itemize}
\tightlist
\item
  HAT 类结构已经达到渐近最优。
\item
  进一步降低存储空间需要新的空间模型。
\end{itemize}

\subsection{4.4 Tarjan 和 Zwick
的改进}\label{tarjan-ux548c-zwick-ux7684ux6539ux8fdb}

Tarjan 和 Zwick 重新设计了空间模型。

他们将空间分为：

\begin{itemize}
\tightlist
\item
  Storage Space：长期保存数组元素的空间。
\item
  Temporary Space：调整过程中使用的临时空间。
\end{itemize}

通过允许更多临时空间，论文突破了传统空间限制。

其结果：

{\def\LTcaptype{none} % do not increment counter
\begin{longtable}[]{@{}ll@{}}
\toprule\noalign{}
项目 & 复杂度 \\
\midrule\noalign{}
\endhead
\bottomrule\noalign{}
\endlastfoot
Storage Space & \(N+O(N^{1/r})\) \\
Temporary Space & \(N+O(N^{1-1/r})\) \\
Access & \(O(1)\) \\
Grow/Shrink & \(O(r)\) 均摊 \\
\end{longtable}
}

参数 \(r\) 控制两者之间的平衡：

\begin{itemize}
\tightlist
\item
  \(r\) 较大时，Storage Space 更接近 \(N\)。
\item
  同时调整操作成本增加。
\end{itemize}

\subsubsection{4.4.1
论文发表后的后续进展}\label{ux8bbaux6587ux53d1ux8868ux540eux7684ux540eux7eedux8fdbux5c55}

Tarjan 和 Zwick 的研究进一步推动了动态数组空间优化方向的发展。

该论文最初以预印本形式公开，随后发表于 SIAM Journal on Computing。\\
论文正式发表后，其提出的 Storage Space 与 Temporary Space
分离思想成为分析动态数组空间效率的重要方法。

\paragraph{（1）工程实现方向}

传统动态数组主要采用倍增策略。

典型实现包括：

\begin{itemize}
\item
  C++ vector。
\item
  Java ArrayList。
\item
  Python list。
\end{itemize}

这些结构具有以下特点：

\begin{itemize}
\item
  实现简单。
\item
  访问效率高。
\item
  CPU 缓存局部性较好。
\end{itemize}

Optimal Resizable Arrays 通过更复杂的数据结构降低空间浪费。\\
这种方法需要维护：

\begin{itemize}
\item
  多级数据块。
\item
  额外索引结构。
\item
  动态调整过程。
\end{itemize}

因此，在实际工程应用中，需要进一步考虑：

\begin{itemize}
\item
  内存分配成本。
\item
  数据访问模式。
\item
  CPU 缓存效率。
\item
  实现维护成本。
\end{itemize}

\paragraph{（2）理论推广方向}

论文提出的核心思想是：

\begin{quote}
通过增加调整过程中的临时空间，降低长期存储空间需求。
\end{quote}

这一思想为其他空间高效数据结构提供了新的设计方向。\\
相关研究方向包括：

\begin{itemize}
\item
  紧凑型动态容器。
\item
  外部存储数据结构。
\item
  内存受限环境中的数据管理。
\end{itemize}

\paragraph{（3）当前影响}\label{ux5f53ux524dux5f71ux54cd}

目前，Optimal Resizable Arrays 主要影响体现在两个方面：

\begin{enumerate}
\def\labelenumi{\arabic{enumi}.}
\item
  理论分析方面。

  该工作提供了新的动态数组空间模型，重新分析了存储空间和调整空间之间的关系。
\item
  数据结构设计方面。

  该工作展示了经典数据结构仍然存在进一步优化空间。虽然其复杂实现限制了大规模工程应用，但它为未来空间高效数据结构设计提供了理论基础。
\end{enumerate}

\subsection{4.5
相关工作总结}\label{ux76f8ux5173ux5de5ux4f5cux603bux7ed3}

{\def\LTcaptype{none} % do not increment counter
\begin{longtable}[]{@{}llll@{}}
\toprule\noalign{}
方法 & 存储空间 & 调整时间 & 特点 \\
\midrule\noalign{}
\endhead
\bottomrule\noalign{}
\endlastfoot
Doubling & \(O(N)\) & \(O(1)\) 均摊 & 简单，工程中常用 \\
Sitarski HAT & \(N+O(\sqrt{N})\) & \(O(1)\) & 降低空间浪费 \\
Brodnik 等结构 & \(N+O(\sqrt{N})\) & \(O(1)\) & 证明空间下界 \\
Tarjan-Zwick & \(N+O(N^{1/r})\) & \(O(r)\) & 利用临时空间突破限制 \\
\end{longtable}
}

Tarjan 和 Zwick
的工作建立在此前动态数组研究基础上，通过重新划分空间模型进一步推进了动态数组优化的理论边界。

\begin{center}\rule{0.5\linewidth}{0.5pt}\end{center}

\section{第5章 下界理论}\label{ux7b2c5ux7ae0-ux4e0bux754cux7406ux8bba}

\subsection{5.1 抽屉原理证明 Brodnik
下界}\label{ux62bdux5c49ux539fux7406ux8bc1ux660e-brodnik-ux4e0bux754c}

Brodnik等人（1999）证明了：\textbf{任何可伸缩数组实现，都必须在某些时刻使用
\(N + \Omega(\sqrt{N})\) 空间}

\textbf{定理 4.1（Brodnik et al.）}：任何支持 \(\text{access}\) 与
\(\text{grow}\) 操作的可伸缩数组，存在某个时刻，已分配的总空间至少为
\(N + \Omega(\sqrt{N})\)

\textbf{证明思路}（抽屉原理）：

考虑连续执行 \(N\) 次 \(\text{grow}\) 操作后的状态。元素被存储在 \(k\)
个内存块中。分两种情况论证：

\begin{enumerate}
\def\labelenumi{\arabic{enumi}.}
\item
  \textbf{\(k \ge \sqrt{N}\)}：每个块都需要一个指针指向该块的起始地址。仅指针数组本身就占用了至少
  \(k \ge \sqrt{N}\) 的空间，而这些指针构成纯粹额外开销
\item
  \textbf{\(k < \sqrt{N}\)}：由抽屉原理，\(N\) 个元素分入少于
  \(\sqrt{N}\) 个块中，必存在某个块的大小超过
  \(N/k > \sqrt{N}\)。该块刚被分配的时刻------即该块为空时，该空块本身构成超过
  \(\sqrt{N}\) 的临时额外空间
\end{enumerate}

两种情况均导致至少 \(\Omega(\sqrt{N})\) 的额外空间，该定理得证

这一定理为后续所有工作设置了基准线------在不区分''常驻空间''与''临时空间''的前提下，\(O(\sqrt{N})\)
的额外空间是绕不开的硬下界

\begin{center}\rule{0.5\linewidth}{0.5pt}\end{center}

Tarjan 与 Zwick 将上述经典下界推广到常驻额外空间 \(s(N)\) 与临时额外空间
\(t(N)\) 之间的乘积下界

\textbf{定理 4.2}：任何 \((s(N), t(N))\)-\(\text{implementation}\)
必须满足 \(s(N) \cdot t(N) \ge N\)（忽略常数因子）

\textbf{完整}：

\begin{enumerate}
\def\labelenumi{\arabic{enumi}.}
\tightlist
\item
  连续执行 \(N\) 次 \(\text{grow}\) 操作。此时 \(N\)
  个元素被存放在某个集合的内存块中，设块的数量为 \(k\)
\item
  每个已分配的块需要一个指针来记录其地址。这些指针存储在索引结构中，占用至少
  \(k\) 个 word。因此指针开销本身已构成额外空间的下界：\(k \le s(N)\)
\item
  由抽屉原理，\(N\) 个元素分布在 \(k\) 个块中，必存在某个块的大小至少为
  \(\lceil N/k \rceil \ge N/s(N)\)
\item
  考虑该最大块的分配时刻。在刚分配完成未填入任何元素时，该块的全部
  \(N/s(N)\) 个 word 均为''临时额外占用''的空间
\item
  由定义，\(\text{grow}\) 过程中的临时额外空间不得超过 \(t(N)\)，因此
  \(t(N) \ge N/s(N)\)
\item
  移项即得 \(s(N) \cdot t(N) \ge N\)。推导完毕
\end{enumerate}

\begin{center}\rule{0.5\linewidth}{0.5pt}\end{center}

\textbf{推论 4.3}：若 \(s(N) = O(N^{1/r})\)，则必有
\(t(N) = \Omega(N^{1-1/r})\)

\textbf{推导}：很简单，将 \(s(N) = O(N^{1/r})\) 代入定理 4.2 即可：

\[t(N) \ge \frac{N}{s(N)} \ge \frac{N}{c \cdot N^{1/r}} = \frac{1}{c} \cdot N^{1-1/r} = \Omega(N^{1-1/r})\]

下面几章证明这个下界是紧的。

\begin{center}\rule{0.5\linewidth}{0.5pt}\end{center}

\section{第6章 上界构造}\label{ux7b2c6ux7ae0-ux4e0aux754cux6784ux9020}

\subsection{6.1 r=3 特例构造}\label{r3-ux7279ux4f8bux6784ux9020}

论文先用 \(r=3\) 的特例来帮助理解

\subsubsection{6.1.1
数据结构与空间分析}\label{ux6570ux636eux7ed3ux6784ux4e0eux7a7aux95f4ux5206ux6790}

设参数 \(B\) 满足 \(N^{1/3} \le B \le 4N^{1/3}\)，于是有：

\begin{itemize}
\tightlist
\item
  \textbf{Large blocks}：大小为
  \(B^2\)，稳态下\textbf{总是满的}，数量至多 \(B\) 个
\item
  \textbf{Small blocks}：大小为 \(B\)，最多 \(2B\)
  个，仅最末一块可部分填充。充当''缓冲区''吸收局部变化，避免频繁触及
  large block 层
\item
  \textbf{索引指针}：合计 \(O(B) = O(N^{1/3})\) 个
\item
  \textbf{稳态额外空间}：\(s(N) = O(N^{1/3})\)（指针 + 至多一个未满
  small block）
\item
  \textbf{临时额外空间}：\(t(N) = O(B^2) = O(N^{2/3})\)（合并/拆分时分配的新块）
\end{itemize}

这个数据结构类似于一个类似 2 位的冗余 B 进制计数器

\subsubsection{6.1.2 核心操作}\label{ux6838ux5fc3ux64cdux4f5c}

\textbf{Access（访问元素）}：

元素在块内外的顺序与逻辑顺序一致，该数据结构寻址等同于二级索引，所以时间复杂度是为
\(O(1)\)

\textbf{Grow（追加元素）}：

\begin{enumerate}
\def\labelenumi{\arabic{enumi}.}
\tightlist
\item
  最末 small block 未满 → 直接写入，\(O(1)\)
\item
  Small block 已满但数量未达 \(2B\) → 分配新 small block 写入，\(O(1)\)
\item
  Small blocks 满 \(2B\) 个 → 触发 \textbf{Combine}：取 \(B\) 个满载
  small block，将其全部元素拷贝到一个新分配的 \(B^2\) large block
  中，释放旧 small block，新 large block 追加到高层，类似于 \(B\)
  进制计数器的\textbf{进位}------低位积累满 \(B\)
  个单元时合并为一个高位单元
\end{enumerate}

\textbf{Shrink（删除末尾元素）}：

\begin{enumerate}
\def\labelenumi{\arabic{enumi}.}
\tightlist
\item
  最末 small block 非空 → 直接删除，\(O(1)\)。
\item
  最末 small block 为空 → 触发 \textbf{Split}：取最末 large
  block，将其末尾 \(B\) 个元素拆分到 \(B\) 个新 small block 中，释放原
  large block，然后从 small block 层删除，类似于 \(B\)
  进制计数器的\textbf{借位}------将一个大额面值拆为小面值
\end{enumerate}

\subsubsection{6.1.3 时间复杂度}\label{ux65f6ux95f4ux590dux6742ux5ea6}

{\def\LTcaptype{none} % do not increment counter
\begin{longtable}[]{@{}lll@{}}
\toprule\noalign{}
操作 & 复杂度 & 说明 \\
\midrule\noalign{}
\endhead
\bottomrule\noalign{}
\endlastfoot
\(\text{access}\) & \(O(1)\) & 两级索引，常数时间定位 \\
\(\text{grow}/\text{shrink}\)（普通） & \(O(1)\) & 直接在小块层操作 \\
\end{longtable}
}

\begin{center}\rule{0.5\linewidth}{0.5pt}\end{center}

\subsection{6.2 一般 r 级构造}\label{ux4e00ux822c-r-ux7ea7ux6784ux9020}

将 r=3 推广到任意整数 \(r \ge 2\)，r=3 模拟的是一个** 2 位的冗余 \(B\)
进制计数器\textbf{， 推广到 \(r \ge 2\) 模拟的是一个} r-1 位的冗余 \(B\)
进制计数器**

\subsubsection{6.2.1
数据结构与空间分析}\label{ux6570ux636eux7ed3ux6784ux4e0eux7a7aux95f4ux5206ux6790-1}

设参数 \(B\) 满足 \(N^{1/r} \le B < 4N^{1/r}\)（为 2 的幂），第 \(i\)
层（\(i = 1, \dots, r-1\)）由大小为 \(B^i\) 的数据块组成。与
§6.1（\(r=3\)）的对比如下：

{\def\LTcaptype{none} % do not increment counter
\begin{longtable}[]{@{}
  >{\raggedright\arraybackslash}p{(\linewidth - 6\tabcolsep) * \real{0.1176}}
  >{\raggedright\arraybackslash}p{(\linewidth - 6\tabcolsep) * \real{0.3725}}
  >{\raggedright\arraybackslash}p{(\linewidth - 6\tabcolsep) * \real{0.3137}}
  >{\raggedright\arraybackslash}p{(\linewidth - 6\tabcolsep) * \real{0.1961}}@{}}
\toprule\noalign{}
\begin{minipage}[b]{\linewidth}\raggedright
\end{minipage} & \begin{minipage}[b]{\linewidth}\raggedright
§6.1（\(r=3\) 特例）
\end{minipage} & \begin{minipage}[b]{\linewidth}\raggedright
§6.2（一般 \(r\)）
\end{minipage} & \begin{minipage}[b]{\linewidth}\raggedright
说明
\end{minipage} \\
\midrule\noalign{}
\endhead
\bottomrule\noalign{}
\endlastfoot
层数 & 2 层（small / large） & \(r-1\) 层 & - \\
块大小 & \(B\)、\(B^2\) & \(B, B^2, \dots, B^{r-1}\) & \(B^i\)
幂次序列 \\
每层容量 & 各至多 \(2B\) 个块 & 各至多 \(2B\) 个块 & 不变，始终是冗余
\(B\) 进制 \\
满载层 & large block（第 2 层） & 第 \(2\) 至 \(r-1\) 层 & 最高层满载 →
除第 1 层外均满载 \\
缓冲区 & small block 层（第 1 层） & 第 1 层（块大小 \(B\)） &
始终是底层的最小块进行增删的缓冲 \\
索引指针 & \(O(B) = O(N^{1/3})\) & \(O(rB) = O(rN^{1/r})\) &
指针数随层数线性增长 \\
\(s(N)\) & \(O(N^{1/3})\) & \(O(rN^{1/r})\) &
\(N^{1/3} \to N^{1/r}\)，多因子 \(r\) \\
\(t(N)\) & \(O(N^{2/3})\) & \(O(N^{1-1/r})\) &
\(N^{2/3} \to N^{1-1/r}\) \\
\(s(N) \cdot t(N)\) & \(O(N)\) & \(O(rN)\) & 均达到定理 4.2 的
\(\Theta(N)\) 下界 \\
\end{longtable}
}

\subsubsection{6.2.2 核心操作}\label{ux6838ux5fc3ux64cdux4f5c-1}

\textbf{Access（访问元素）}：

元素在块内外的顺序与逻辑顺序一致，通过位运算（msb、移位等）在 \(O(1)\)
时间内定位下标 \(i\) 所属的层和块内偏移。

\textbf{Grow（追加元素）}：

\begin{enumerate}
\def\labelenumi{\arabic{enumi}.}
\tightlist
\item
  第 1 层有空位 → 直接写入最末块，\(O(1)\)
\item
  第 1 层满但未达 \(2B\) → 分配新 \(B\) 大小块写入，\(O(1)\)
\item
  第 1 层满 \(2B\) 块 → 触发 \textbf{Combine}：找首个有空位的高层
  \(k\)（\(n_k < 2B\)），对 \(i = k-1\) 到 1，将第 \(i\) 层 \(B\)
  个满块合并为一个 \(B^{i+1}\) 大小的块，释放旧块，新块提升到第 \(i+1\)
  层。类似于 \(B\) 进制计数器的\textbf{进位}------低位积累满 \(B\)
  个单元时合并为一个高位单元
\end{enumerate}

\textbf{Shrink（删除末尾元素）}：

\begin{enumerate}
\def\labelenumi{\arabic{enumi}.}
\tightlist
\item
  第 1 层有元素 → 直接删除末尾元素，\(O(1)\)
\item
  第 1 层无块 → 触发 \textbf{Split}：找最底层有块的高层
  \(k\)（\(n_k > 0\)），将其一个 \(B^k\) 块拆为 \(B\) 个 \(B^{k-1}\)
  块，逐级下传至第 1 层，然后删除。类似于 \(B\)
  进制计数器的\textbf{借位}------将一个大额面值拆为小面值
\end{enumerate}

\textbf{Rebuild（全局重建）}：

当 \(N\) 超出 \([B^r/8, B^r]\) 时，\(B\)
加倍或减半，逐块搬运重建。每次仅分配一个 \(B^{r-1}\)
新块，搬空旧块立即释放，确保重建中空间始终不超过 \(N + O(N^{1-1/r})\)。

\subsubsection{6.2.3 均摊分析（Credit
论证）}\label{ux5747ux644aux5206ux6790credit-ux8bbaux8bc1}

核心思想：考虑未来可能发生的拷贝代价提前分配 credit

\textbf{Grow 方向}：每个新元素分配 \(r-1\) 个 credit 进入第 1
层。每上升一层消耗 1 个 credit 支付拷贝代价。credit 始终非负------第 1
层最多（\(r-1\)），第 \(r-1\) 层最少（\(1\)）\textbf{不变量}：元素
credit 数 = \(r - \text{所在层数}\)。均摊代价 = \(O(1)\) + \(r-1\) =
\(O(r)\)

\textbf{Shrink 方向}：每次 shrink 向每层分配 \(O(1)\) credit。层 \(k\)
的 Split 拷贝代价 \(O(B^k)\) 由该层积累的 credit 支付；两次同层 Split
之间至少经历 \(B^k\) 次 shrink，足够涵盖拷贝代价

\textbf{Rebuild}：代价 \(O(N)\)，间隔 \(\Omega(N)\) 次操作，摊还
\(O(1)\)

\subsubsection{6.2.4 时间复杂度}\label{ux65f6ux95f4ux590dux6742ux5ea6-1}

{\def\LTcaptype{none} % do not increment counter
\begin{longtable}[]{@{}
  >{\raggedright\arraybackslash}p{(\linewidth - 4\tabcolsep) * \real{0.3000}}
  >{\raggedright\arraybackslash}p{(\linewidth - 4\tabcolsep) * \real{0.4000}}
  >{\raggedright\arraybackslash}p{(\linewidth - 4\tabcolsep) * \real{0.3000}}@{}}
\toprule\noalign{}
\begin{minipage}[b]{\linewidth}\raggedright
操作
\end{minipage} & \begin{minipage}[b]{\linewidth}\raggedright
复杂度
\end{minipage} & \begin{minipage}[b]{\linewidth}\raggedright
说明
\end{minipage} \\
\midrule\noalign{}
\endhead
\bottomrule\noalign{}
\endlastfoot
\(\text{access}\) & \(O(1)\) & 位运算定位，常数条机器指令 \\
\(\text{grow}/\text{shrink}\)（普通） & \(O(1)\) & 直接在第 1 层操作 \\
\(\text{Combine}/\text{Split}\) & 均摊 \(O(1)\) & 进位/借位代价由 credit
覆盖 \\
\(\text{Rebuild}\) & 均摊 \(O(1)\) & 代价 \(O(N)\)，每 \(\Omega(N)\)
次操作触发 \\
\end{longtable}
}

\begin{center}\rule{0.5\linewidth}{0.5pt}\end{center}

\subsection{6.3
空间优化变换}\label{ux7a7aux95f4ux4f18ux5316ux53d8ux6362}

§6.2 给出 \(s(N) = O(rN^{1/r})\) 带了一个因子 \(r\)。当 \(r\)
为固定常数时无所谓，仍为最优实现，但如果 r 随着 N 增长（例如
\(r = \log \log N\)），这个 r
就不能简单地直接忽略。于是本节证明额外空间量可以降至
\(s(N) = O(N^{1/r})\)

\subsubsection{6.3.1
变换针对什么时机？}\label{ux53d8ux6362ux9488ux5bf9ux4ec0ux4e48ux65f6ux673a}

§6.2 的数据结构是标准的：每次 Combine 会分配一个 \(B^{i+1}\)
大小的新块，把 \(B\) 个 \(B^i\)
小块中的元素挨个拷贝进去，然后释放那些小块。在''新块刚分配好、旧块还没释放''的瞬间，临时额外空间达到了峰值------这就是
\(t(N) = O(N^{1-1/r})\)。

变换的时机就是在这个分配的瞬间。取
\(T = N^{1-1/r}\)（即临时空间峰值的量级）。每当结构要向系统申请一个大小为
\(b\) 的新块时：

\begin{itemize}
\tightlist
\item
  如果 \(b \le T\)，正常分配，一个块搞定。
\item
  如果 \(b > T\)，不分配一块大的，改为分配 \(\lceil b / T \rceil\)
  个大小为 \(T\)
  的块------相当于把那个大块切成了一串小方块。元素照旧按顺序填进去，读完一块接着下一块。
\end{itemize}

功能完全等价：拷贝总量不变，元素顺序不变，access 多一步块间跳转但仍然是
\(O(1)\)。

\subsubsection{6.3.2
为什么指针反而变少了？}\label{ux4e3aux4ec0ux4e48ux6307ux9488ux53cdux800cux53d8ux5c11ux4e86}

关键后果：变换之后，\textbf{系统里没有任何块的尺寸超过
\(T\)}。所有已分配空间------数据块、索引块、元数据------每个块的大小 ≤
\(T\)。

原来 §6.2 的指针开销是分层的：\(r-1\) 层索引，每层各自维护至多 \(2B\)
个指向数据块的指针，层级叠加出 \(r\)
因子。变换之后分层没意义了------所有块都是 \(T\)
大小，一层统一管理即可。总块数（也就是总指针数）不超过总分配空间除以
\(T\)：

\[\text{总块数} \le \frac{N + s(N) + t(N)}{T} \approx \frac{N}{N^{1-1/r}} = N^{1/r}\]

\subsubsection{6.3.3 形式化}\label{ux5f62ux5f0fux5316}

论文将上述操作归纳为一个通用引理，然后将它作用到 §6.2 构造的一个变体上：

\begin{quote}
\textbf{引理 7.2}：设有标准 \((s(N), O(N))\)-实现。对任意
\(t(N)\)，可将其所有''分配新块''操作中的大块替换为 \(t(N)\)
大小的子块，得到
\(s'(N) = s(N) + O(N/t(N))\)、\(t'(N) = s(N) + O(t(N))\)
的新实现，时间界不变。
\end{quote}

\begin{quote}
\textbf{定理 7.3}：取 §6.2 中参数为 \(2r\)
的构造（\(s = O(rN^{1/(2r)}), t = O(N^{1-1/(2r)})\)），对其施加引理
7.2（取 \(t(N) = N^{1-1/r}\)），得
\(s'(N) = O(rN^{1/(2r)} + N^{1/r})\)。在
\(r \le \frac{1}{2}\frac{\log N}{\log \log N}\)
条件下，\(rN^{1/(2r)} \le N^{1/r}\)，故
\(s'(N) = O(N^{1/r})\)、\(t'(N) = O(N^{1-1/r})\)。
\end{quote}

\subsubsection{6.3.4 实践意义}\label{ux5b9eux8df5ux610fux4e49}

\(r\) 随 \(N\) 增长的场景：例如 \(r = \Theta(\log N)\) 时
\(rN^{1/r} = \Theta(\log N)\)，变换后降到
\(N^{1/r} = \Theta(1)\)------把额外空间压到了常数级。

\begin{center}\rule{0.5\linewidth}{0.5pt}\end{center}

\section{第7章
紧致性证明}\label{ux7b2c7ux7ae0-ux7d27ux81f4ux6027ux8bc1ux660e}

第6章证明了 \(O(r)\)
均摊的\textbf{上界}。本章使用一个抽象博弈证明无法更好了------其推论直接给出
\(\Omega(r)\) 的下界。

\subsection{7.1 Growth Game
模型与精确解}\label{growth-game-ux6a21ux578bux4e0eux7cbeux786eux89e3}

\subsubsection{7.1.1 博弈定义}\label{ux535aux5f08ux5b9aux4e49}

博弈设定：\(N\) 个元素逐个插入，玩家维护 \(k\) 个子数组
\(A_1, \dots, A_k\)（初始为空），最多允许 \(\ell\)
个空位。空子数组总是在最前面，第一个非空子数组至多 \(\ell\)
个空位，其余全满。插入时分三种情况：

{\def\LTcaptype{none} % do not increment counter
\begin{longtable}[]{@{}
  >{\raggedright\arraybackslash}p{(\linewidth - 4\tabcolsep) * \real{0.3333}}
  >{\raggedright\arraybackslash}p{(\linewidth - 4\tabcolsep) * \real{0.3333}}
  >{\raggedright\arraybackslash}p{(\linewidth - 4\tabcolsep) * \real{0.3333}}@{}}
\toprule\noalign{}
\begin{minipage}[b]{\linewidth}\raggedright
情况
\end{minipage} & \begin{minipage}[b]{\linewidth}\raggedright
操作
\end{minipage} & \begin{minipage}[b]{\linewidth}\raggedright
代价
\end{minipage} \\
\midrule\noalign{}
\endhead
\bottomrule\noalign{}
\endlastfoot
(a) 有空位 & 直接写入 & \(1\) \\
(b) 有空子数组 & 分配大小 \(\ell+1\) 的新子数组，写入 & \(1\) \\
(c) 全部 \(k\) 个满 & \textbf{选 \(i \in [k]\)}，合并
\(A_i, \dots, A_1\) 到新块，写入 & \(1 + \sum_{j=1}^i a_j\) \\
\end{longtable}
}

(a)(b) 无选择。博弈的核心在 (c)：\(k\) 种合法移动------选 \(i=1\)
代价小但碎片多，选 \(i=k\) 代价大但干净。目标：使总代价 \(C_{N,k,\ell}\)
最小。

\subsubsection{7.1.2 博弈分析}\label{ux535aux5f08ux5206ux6790}

\textbf{第一步：消除 \(\ell\)。} 每次分配新子数组留下 \(\ell\)
个空位后，接下来 \(\ell\) 次插入必然直接填空，不需要决策。每 \(\ell+1\)
次操作中仅 1 次是真正的决策。将 \(\ell+1\)
个元素打包为一个逻辑单元，博弈退化为无空位版本：

\[C_{N,k,\ell} = (\ell+1) \cdot C_{N/(\ell+1),\, k,\, 0}, \qquad A_{N,k,\ell} = A_{N/(\ell+1),\, k,\, 0}\]

此后只讨论 \(\ell = 0\)，记 \(C_{N,k} = C_{N,k,0}\)。

\textbf{第二步：建立递推}。 状态用向量
\(\mathbf{a} = (a_1, \dots, a_k)\)
描述（\(\sum a_i = N\)，空子数组在前缀）。令 \(C(\mathbf{a})\)
为从全零到达 \(\mathbf{a}\) 的最小代价。到达 \(\mathbf{a}\)
的最后一步必定是最后一次改变 \(A_k\) 的合并，它将系统带入
\((0, \dots, 0, a_k)\) 态，此后 \(A_k\) 不再变化。由此（引理 8.4）：

\[C_k(a_1, \dots, a_k) = C_k(0, \dots, 0, a_k) + C_{k-1}(a_1, \dots, a_{k-1}) \tag{1}\]

\[C_k(0, \dots, 0, a_k) = C_{a_k-1,\, k} + a_k \tag{2}\]

\begin{enumerate}
\def\labelenumi{(\arabic{enumi})}
\tightlist
\item
  将代价拆为''到达 \(a_k\)``+''前 \(k-1\) 位独立操作''两部分。(2)
  刻画最后一步：合并前 \(a_k-1\) 个元素所在的全部子数组，代价
  \(a_k\)，合并前的代价为
  \(C_{a_k-1,k}\)。反复展开得到独立求和式，将任意状态的代价关联到更小参数的最优代价------归纳证明的骨架。
\end{enumerate}

\textbf{第三步：闭式解。} 通过对 \((N, k)\) 双层归纳求解递推。令 \(n\)
满足 \(\binom{n+k-1}{k} - 1 \le N \le \binom{n+k}{k} - 1\)，定理 8.6
给出：

\[C_{N,k} = (N+1)n - \binom{n+k}{k+1}, \qquad A_{N,k} = \frac{C_{N,k}}{N} \ge \frac{n-1}{2}\]

最优状态下子数组大小被二项式系数夹逼：\(\binom{n+i-2}{i} \le a_i \le \binom{n+i-1}{i}\)------``大块多大、小块多小''由组合数决定，不是随意定的。证明核心是扰动论证：若
\(a_i\) 偏离上述范围，挪一个元素即可严格降低总代价。

\textbf{第四步：构造策略。} 最优代价可以达到的。策略维护二项式计数器
\((b_1, \dots, b_k)\)，子数组大小
\(a_i = \binom{b_i + i - 1}{i}\)。每次找最小的 \(i\) 使
\(b_i < b_{i+1}\)，合并 \(A_1, \dots, A_i\)，\(b_i\)
加一、低位归零。这可以看作 §6.2 中 \(B\)
进制计数器的理论最优版------进位基数不固定，由二项式系数动态决定（§6.2
的 \(2B\) 冗余是为同时支持 grow/shrink 的妥协；这里只有
grow，可达精确最优）。

\begin{center}\rule{0.5\linewidth}{0.5pt}\end{center}

\subsection{\texorpdfstring{7.2 归约至数据结构 \(\Omega(r)\)
下界}{7.2 归约至数据结构 \textbackslash Omega(r) 下界}}\label{ux5f52ux7ea6ux81f3ux6570ux636eux7ed3ux6784-omegar-ux4e0bux754c}

归约分四步：

\textbf{第一步：块 → 子数组。}
数据结构中维护的每个内存块对应博弈中的一个子数组 \(A_i\)，块中元素数即
\(a_i\)。当数据结构触发 Combine 合并若干块时，对应博弈中的 (c)
类合并操作，元素拷贝量恰好等于博弈中的移动代价。

\textbf{第二步：额外空间 → \(k\) 和 \(\ell\)。} 常驻额外空间
\(s(N) = O(N^{1/r})\) 意味着块的数目受限于
\(O(N^{1/r})\)。更精细地分析，多层索引结构使总块数
\(k \approx r \cdot N^{1/r}\)。部分填充的块对应空位
\(\ell\)，二者量级相当。

\textbf{第三步：映射到博弈下界。} \(N\) 次 grow 操作自然诱导一个
\((N, k, \ell)\)-增长博弈的合法策略。由归约 \(\ell \to 0\) 及推论
8.13，该博弈的均摊代价满足 \(A_{N,k} \ge (n-1)/2\)。当
\(k = \Theta(rN^{1/r})\) 时，由二项式系数的渐近性质可得
\(n = \Theta(r)\)，故 \(A_{N,k} = \Omega(r)\)。

\textbf{第四步：结论。}
博弈的均摊代价是数据结构真实拷贝次数的下界，因此：

\begin{quote}
\textbf{定理 9.1}：当 \(r=O(\log N)\) 时，任何使用 \(N + O(rN^{1/r})\)
存储空间的标准实现，其 grow 操作的均摊代价为 \(\Omega(r)\)。
\end{quote}

完整归约链条为：

\begin{center}
\begin{tabular}{@{}l@{}}
标准数据结构（\(N + O(rN^{1/r})\) 存储）\\
\(\quad\downarrow\) 块 \(\to\) 子数组，合并操作 \(\to\) 博弈移动\\
\(\quad\downarrow\)\\
\((N, k, \ell)\)-增长博弈（\(k \approx rN^{1/r}\)，\(\ell \approx N^{1/r}\)）\\
\(\quad\downarrow\) 引理 8.2：吸收空位\\
\(\quad\downarrow\)\\
\((N/(\ell+1), k)\)-增长博弈（\(\ell = 0\)）\\
\(\quad\downarrow\) 定理 8.6 + 推论 8.13：\(A_{N,k} = \Omega(r)\)\\
\(\quad\downarrow\)\\
Grow 均摊代价 \(= \Omega(r)\)
\end{tabular}
\end{center}

结合第 6 章的上界 \(O(r)\)，论文在 grow/shrink
的均摊时间上达到了紧致（tight）。下表汇总全部维度的上下界对比：

{\def\LTcaptype{none} % do not increment counter
\begin{longtable}[]{@{}
  >{\raggedright\arraybackslash}p{(\linewidth - 6\tabcolsep) * \real{0.1765}}
  >{\raggedright\arraybackslash}p{(\linewidth - 6\tabcolsep) * \real{0.4706}}
  >{\raggedright\arraybackslash}p{(\linewidth - 6\tabcolsep) * \real{0.1765}}
  >{\raggedright\arraybackslash}p{(\linewidth - 6\tabcolsep) * \real{0.1765}}@{}}
\toprule\noalign{}
\begin{minipage}[b]{\linewidth}\raggedright
维度
\end{minipage} & \begin{minipage}[b]{\linewidth}\raggedright
上界（第 6 章）
\end{minipage} & \begin{minipage}[b]{\linewidth}\raggedright
下界
\end{minipage} & \begin{minipage}[b]{\linewidth}\raggedright
结论
\end{minipage} \\
\midrule\noalign{}
\endhead
\bottomrule\noalign{}
\endlastfoot
稳态存储 & \(N + O(N^{1/r})\) & \(\ge N\)（显然） & 渐近最优 \\
临时 resize 空间 & \(N + O(N^{1-1/r})\) &
\(N + \Omega(N^{1-1/r})\)（定理 4.2） & \textbf{紧致} \\
Grow/Shrink 时间 & \(O(r)\) 均摊 & \(\Omega(r)\) 均摊（定理 9.1） &
\textbf{紧致} \\
Access 时间 & \(O(1)\) 最坏 & \(\Omega(1)\)（显然） & 最优 \\
\end{longtable}
}

论文在存储空间、临时空间、更新时间三个维度上同时达到渐近最优------使''Optimal
Resizable Arrays''名副其实

\begin{center}\rule{0.5\linewidth}{0.5pt}\end{center}

\section{第8章
讨论与评价}\label{ux7b2c8ux7ae0-ux8ba8ux8bbaux4e0eux8bc4ux4ef7}

通读论文其中最重要的几个关键设计：

\textbf{（1）把存储空间和临时空间分开。} Brodnik
的下界说''不管什么时候，额外空间至少
\(\Omega(\sqrt{N})\)``。但实际场景中，一个程序有上百个数组，只有正在
resize
的那个才需要临时空间，其余都不需要。把两者分开考虑就可以绕开了这个下界------定理
4.2 的 \(s(N) \cdot t(N) \ge N\) ：二者乘积受限，但各自有可调节余量

\textbf{（2）用 \(B\) 进制计数器来组织块。} 两级结构（\(B\) 和
\(B^2\)）能把额外空间压到
\(O(\sqrt{N})\)。多叠几层------\(B, B^2, \dots, B^{r-1}\)便可把空间继续下压，每
\(B\)
个小块合并成一个更大的块，与进位逻辑相同。层数越多空间越省，但每个元素最多被搬
\(r-1\) 次------得到了 \(O(r)\)

\textbf{（3）每层允许 \(2B\) 个块而不只是 \(B-1\)。} \(B-1\) 是标准
\(B\) 进制计数器的容量，但那样只能 grow 不能
shrink。多留一倍的缓冲空间，Shrink
拆下来的小块有地方放，不会刚合并完又拆开。所以使用冗余 \(B\)
进制计数器进行双向缓冲

\textbf{（4）用 credit 进行均摊分析。} Grow 的 credit
在元素层面------每升一层花 1 个 credit，不变量自已就是证明。Shrink 的
credit 在层的层面------每层靠 shrink 的''存款''来支付 Split。Rebuild
再额外存一份。三方的credit各自独立，加起来得到了 \(O(r)\)。

\begin{center}\rule{0.5\linewidth}{0.5pt}\end{center}

\section{第9章
局限性与开放问题}\label{ux7b2c9ux7ae0-ux5c40ux9650ux6027ux4e0eux5f00ux653eux95eeux9898}

\subsection{9.1
下界的适用范围}\label{ux4e0bux754cux7684ux9002ux7528ux8303ux56f4}

论文证明了 grow 操作的均摊时间下界为
\(\Omega(r)\)，但该下界只适用于论文定义的 standard implementation。

也就是说，论文尚未完全排除非标准实现获得更好结果的可能。作者认为，若对元素和指针加入合适的不可压缩性假设，非标准实现也不会有渐近优势。

因此，一个开放问题是：如何把 \(\Omega(r)\)
的下界推广到更一般的算法模型。

\subsection{9.2 Growth Game
证明较复杂}\label{growth-game-ux8bc1ux660eux8f83ux590dux6742}

为了证明下界，论文构造并精确求解了 Growth
Game。该博弈描述了连续插入时如何合并数据块，才能使总复制成本最小。

精确解需要较复杂的组合分析。论文指出，若只需要渐近意义上的 \(\Omega(r)\)
下界，可能存在更简单的归纳或势能证明方法。

\subsection{9.3
工程实现的限制}\label{ux5de5ux7a0bux5b9eux73b0ux7684ux9650ux5236}

该结构需要维护多级数据块、索引结构、Combine、Split
和重建过程，因此实现比传统倍增数组复杂。

参数 \(r\) 也需要在空间和时间之间取舍：\(r\) 较大时，常驻存储空间更接近
\(N\)，但 grow 和 shrink 的均摊时间会增加到
\(O(r)\)。未来可进一步比较该结构与传统倍增数组、HAT
等结构在真实内存分配器和缓存环境下的表现。

\begin{center}\rule{0.5\linewidth}{0.5pt}\end{center}

\section{第10章 总结}\label{ux7b2c10ux7ae0-ux603bux7ed3}

Tarjan 和 Zwick
研究了动态数组的空间与时间权衡问题。传统倍增数组访问快、扩容均摊时间低，但会保留较多空闲空间。

论文的核心改进是区分 Storage Space 和 Temporary
Space：允许扩容或缩容时使用更多临时空间，从而减少稳定状态下的长期存储空间。

论文证明，对于参数 \(r \geq 2\)，可以实现：

\[
\text{常驻存储空间 } N + O(N^{1/r}),\qquad
\text{临时空间 } N + O(N^{1-1/r}),\qquad
\text{grow/shrink 均摊时间 } O(r).
\]

同时，随机访问保持为 \(O(1)\) 最坏时间。Growth Game 又给出 grow 操作的
\(\Omega(r)\)
均摊时间下界，因此该结果在论文的标准模型中达到了紧致的空间---时间权衡。

该工作说明：经典数据结构仍然可以通过重新划分资源模型获得新的优化结果。虽然结构实现较复杂，但它为未来空间高效数据结构的设计提供了理论基础。

\end{document}
